\documentclass[11pt]{article}
\usepackage{mathblatt}
% Präambel des Setzers (werkzeuge/setzer-praeambel.tex, wortgleich); dazu \gteil (Muster T2U)
\makeatletter
\setlength{\mbpfbe}{0mm}% keine Punkte (pruefung.md § 5: nur beim Prüfstein)
\newlength{\gEin}\setlength{\gEin}{0pt}
\let\gSchrift\relax
\renewcommand{\pfkreuzzeile}[1]{\par\noindent\hangindent3.2em\hangafter1\hspace*{1.6em}$\square$\ #1\par}
\newcommand{\gkreuz}[1]{$\square$\ #1\hspace{2em}}
\newcommand{\gkreuzzeile}[2]{\par\noindent\hangindent2.2em\hangafter1\mbox{#1\,$\square$}\ #2\par}% Buchstabe, Option: eine je Zeile (neben Skizze)
% links, was man ansieht, rechts, was man tut; Spaltengrenze überall an derselben Stelle (gemeinsam.md „Satz“)
% v0.17: Fußeintrag „Nr: Wert“ – kurz in einer Box (nie getrennt), lang umbrechbar mit der Nummer am Anfang
% (vorher \mbox: lange Lösungen liefen über den Rand, AS40 138 pt); robust, weil \fusshilfe in die .aux schreibt
\NewDocumentCommand{\gfuss}{m m}{\sbox0{#1:~#2}\ifdim\wd0<0.45\textwidth\mbox{#1:~#2}\else\mbox{#1:}\nobreak\ #2\fi}
\newsavebox{\gbox}\newlength{\gLinks}
% v0.10 (Lehrer 13.10.): keine starre Halbierung – die linke Spalte ist so breit wie die Figur (+ 6 mm), die
% Antwort steht direkt daneben; ist die Figur breiter als 0,62 Zeile, steht die Antwort darunter
\newcommand{\gzwei}[2]{\par\smallskip\noindent\sbox{\gbox}{#1}%
  \setlength{\gLinks}{\dimexpr\wd\gbox+6mm\relax}%
  \ifdim\gLinks>0.62\linewidth
    \usebox{\gbox}\par\smallskip\noindent\begin{minipage}{\linewidth}#2\end{minipage}%
  \else
    \begin{minipage}[t]{\gLinks}\vspace{0pt}\usebox{\gbox}\end{minipage}%
    \begin{minipage}[t]{\dimexpr\linewidth-\gLinks\relax}\vspace{0pt}\raggedright #2\end{minipage}%
  \fi\par}
% v0.9 (12.10., pruefung.md § 5 „Marke mit B, Bedingungen rechts“): drittes Argument von pfaufg = Bedingungen
% der Prüfung (v0.10: \pfohnetr = TR durchgestrichen, „FOR“), klein grau im rechten Rand auf der Grundlinie der Nummer, Zeilen mit \\
\newcommand{\gbedrand}[2]{\ifblank{#2}{}{\rlap{\hspace*{\dimexpr#1+2mm\relax}%
  \parbox[t]{15mm}{\raggedright\scriptsize\leavevmode\color{mbgrau}\strut#2}}}}
% Teil eines Bündels (a, b …): Marke links im Rand, Bedingungen rechts im Rand, auf der Zeile des Teils
\newcommand{\gteilrand}[2]{\llap{\mbpfrandmarke{#1}\hspace*{\mbpfnr}}\gbedrand{\linewidth}{#2}}
% Teil eines Bündels: Buchstabe in eigener Spalte, Inhalt eingerückt (\gEin = Einzug, damit die Spaltengrenze von
% \gzwei an derselben Stelle bleibt)
\newenvironment{gbteil}[3]{\par\smallskip\noindent\gteilrand{#2}{#3}\makebox[\mbpfteil][l]{\textbf{#1}}%
  \begin{minipage}[t]{\dimexpr\linewidth-\mbpfteil\relax}\setlength{\gEin}{\mbpfteil}\raggedright\strut\ignorespaces}%
 {\end{minipage}\par\addvspace{6pt}}% v0.18 (PG80 Nr. 1): Abstand, damit die Linie von a) nicht an b) rutscht
% Grafik höchstens so breit wie die Spalte (Bündel nebeneinander, Koordinatensystem ohne Karo)
\newsavebox{\gpbox}
\newcommand{\gpassend}[1]{\sbox{\gpbox}{#1}\ifdim\wd\gpbox>\linewidth\typeout{GPASSEND: verkleinert \the\wd\gpbox\space auf \the\linewidth}\resizebox{\linewidth}{!}{\usebox{\gpbox}}\else\usebox{\gpbox}\fi}
% Skizzenfeld (v0.10, Lehrer 13.10.): Karo ab dem linken Rand über die Zeile, „Skizze“ klein grau oben links im Feld
% (wie \pfskizzenfeld der Vorlage); der rechte Rand bleibt den Bedingungen
\newcommand{\gskizze}[1]{\par\smallskip\noindent\begin{tikzpicture}
  \clip (0,0) rectangle ({\linewidth-1mm},{#1*10mm});
  \draw[black!16,line width=0.3pt,step=5mm] (0,0) grid ({\linewidth},{#1*10mm});
  \node[anchor=north west,inner sep=2pt,font=\scriptsize,text=mbgrau,fill=white] at (0,{#1*10mm}) {Skizze};
  \end{tikzpicture}\par}
\RenewDocumentEnvironment{pfaufg}{m m m}{%
  \begin{lrbox}{\mbaufgabebox}\begin{minipage}[t]{\linewidth}%
  \gSchrift\noindent\hspace*{\gEin}\mbpfrandmarke{#1}\gbedrand{\linewidth-\gEin}{#3}%
  \begin{minipage}[t]{\mbpfnr}\strut\textbf{#2}\end{minipage}%
  \begin{minipage}[t]{\dimexpr\linewidth-\gEin-\mbpfnr\relax}\raggedright\strut\ignorespaces}%
 {\end{minipage}%
  \end{minipage}\end{lrbox}%
  \setlength{\mbaufgabehoehe}{\dimexpr\ht\mbaufgabebox+\dp\mbaufgabebox\relax}%
  \ifdefempty{\gkopfnach}{\par\addvspace{5pt}\Needspace*{\mbaufgabehoehe}}{%
    % Stufenkopf, der mit dieser (ersten) Aufgabe zusammenbleibt: Platz für beide verlangen, dann den Kopf setzen
    % (12.10., Quadratische: ein Kopf stand allein unten, die Aufgabe mit Koordinatensystem auf der nächsten Seite)
    % v0.17 (AS40 „Terme deuten“ allein unten): hinter dem Kopf steht das Sprungziel (\special, ein Knoten, nach
    % dem die Seite brechen durfte) – darum \nobreak nach dem Kopf; Vorhalt 56 pt = Kopf mit Linie und Abständen
    \par\Needspace*{\dimexpr\mbaufgabehoehe+56pt\relax}\gkopfnach\par\nobreak\gdef\gkopfnach{}\addvspace{5pt}}%
  \noindent\usebox{\mbaufgabebox}\par\addvspace{5pt}}
% v0.18 (plan.md § 10, 15.10. l): Abschnittsbuchstabe A, B, C vor dem Namen im Kopf – als Vorsatz, damit der
% Name im TeX gleich bleibt (Bestellkarte, Sprungziele lesen ihn)
\def\gbstakt{}\newcommand{\gbstvor}[1]{\gdef\gbstakt{#1}}
\NewCommandCopy\gpfstufekopfalt\pfstufekopf
\RenewDocumentCommand{\pfstufekopf}{o m m}{\IfValueTF{#1}{\gpfstufekopfalt[#1]}{\gpfstufekopfalt}%
  {\ifx\gbstakt\empty\else\gbstakt\enspace\fi #2}{#3}\gdef\gbstakt{}}
% v0.18 (16.10. e): Begriff grau, ohne Nummer, ohne Rahmen über der ersten Aufgabe der Stufe
\newcommand{\gbegriff}[1]{\par\nobreak\noindent\hspace*{\mbpfnr}\parbox[t]{\dimexpr\linewidth-\mbpfnr\relax}{\raggedright\small\color{mbgrau}#1}\par\nobreak\addvspace{3pt}}
% v0.18 (16.10. g): Lösungsblatt – Kopf mit kräftigem Strich, Abschnitt klein grau als Zwischenzeile
\newcommand{\glkopf}[1]{\noindent{\large\bfseries Lösungen \textperiodcentered{} #1}\par\vspace{3pt}\hrule height 1.2pt\par\medskip}
\newcommand{\glzwischen}[1]{\par\addvspace{8pt}\Needspace*{4\baselineskip}\noindent{\footnotesize\color{mbgrau}#1}\par\nobreak\vspace{1pt}}
% v0.18 (15.10. o, w): Blattende – Kern ohne Rahmen, dann graue Zeile „Mehr: n Aufgaben“
\newcommand{\gkern}[1]{\par\addvspace{14pt}\Needspace*{6\baselineskip}\noindent{\bfseries Kern}\par\nobreak\smallskip
  {\small\raggedright #1\par}}
\newcommand{\gmehr}[1]{\par\addvspace{10pt}\noindent{\footnotesize\color{mbgrau}Mehr: #1 Aufgaben}\par}
% Stufenkopf vormerken: \gkopfvor{<Kopf>} – die nächste pfaufg setzt ihn, sobald beide zusammen Platz haben
\def\gkopfnach{}
\newcommand{\gkopfvor}[1]{\gdef\gkopfnach{#1}}
% Teilaufgabe mit Skizze links in gzwei: Buchstabe in eigener Spalte, darunter Text und Skizze
\newcommand{\gteil}[2]{\makebox[\mbpfteil][l]{\textbf{#1}}\parbox[t]{\dimexpr0.5\textwidth-\gEin-\mbpfnr-5mm-\mbpfteil\relax}{\raggedright #2}}
% Stufenkopf der Vorlage (mathblatt.sty 2026-10-12: Linie, Abstand, 13 pt, \pfjahre) ohne ihren Vorhalt von
% 14 Zeilen: der Setzer fordert vor dem Kopf 0,16 Seitenhöhe (Kopf + erste Aufgabe), mehr nicht
\setlength{\mbpfstufevorhalt}{0pt}
\makeatother
\newcommand{\karo}[1]{\par\smallskip\noindent\begin{tikzpicture}
  \clip (0,0) rectangle ({\linewidth-1mm},{#1*10mm});
  \draw[black!16,line width=0.3pt,step=5mm] (0,0) grid ({\linewidth},{#1*10mm});
  \end{tikzpicture}\par}
\newcommand{\kk}[1]{$\square$\,#1\hspace{0.9em}}
\newcommand{\linie}[1]{\,{\color{black!45}\rule[-1pt]{#1}{0.4pt}}\,}
\newcommand{\pkt}{\textperiodcentered{}}
% v0.14 (14.10.): Stufenkopf mit fünf Punkten (Option = Jahre für die Bestellkarte) oder „zuletzt 19“
\NewDocumentCommand{\gpunkte}{o m}{#2}
\newcommand{\gzuletzt}[1]{zuletzt #1}
% v0.7 (12.10.): Schülerfrage unter dem Stufenkopf; Ankreuzoption als Bild (Buchstabe und Kästchen darunter)
\newcommand{\gfrage}[1]{\par\nobreak\noindent\hspace*{\mbpfnr}{\small\itshape #1}\par\nobreak\addvspace{2pt}}
\newlength{\gbildbreite}\setlength{\gbildbreite}{0.32\linewidth}% Setzer setzt sie je Reihe (vier Bilder: 0,24)
\newcommand{\gbildoption}[3]{\begin{minipage}[t][#3][b]{\gbildbreite}\centering #2\par\smallskip\mbox{#1\,$\square$}\end{minipage}}
% Lineare (12.10.): Bild mit Buchstabe ohne Kästchen (Diagramm zuordnen); Aussage mit „wahr / falsch“ zum Ankreuzen
\newcommand{\gbildname}[3]{\begin{minipage}[t][#3][b]{\gbildbreite}\centering #2\par\smallskip\mbox{\textbf{#1}}\end{minipage}}
\newcommand{\gwahrfalsch}[2]{\par\noindent\hangindent2.2em\hangafter1\makebox[2.2em][l]{(#1)}#2\hfill\mbox{wahr $\square$\quad falsch $\square$}\par}
\begin{document}

\pfheftstil
\pflaufkopf[Abitur GK \pkt{} Prüfung]{Bedingte Wahrscheinlichkeit}{AS40}\pfseitevon
\pfheftkopf{Bedingte Wahrscheinlichkeit}{}
\pfstufenzeile{bedingte Wahrscheinlichkeit gesucht \pkt{} bedingte Wahrscheinlichkeit mit Parameter \pkt{} Terme deuten}

% ===================== Stufe: bedingte Wahrscheinlichkeit gesucht =====================
\gkopfvor{\pfstufekopf{bedingte Wahrscheinlichkeit gesucht}{\gpunkte[23 22]{\textcolor{black!25}{$\circ$}\,\textcolor{black!25}{$\circ$}\,\textcolor{black!25}{$\circ$}\,\textcolor{black!80}{$\bullet$}\,\textcolor{black!80}{$\bullet$}}}\special{pdf:dest (stufe-bedingte-wahrscheinlichkeit-gesucht) [@thispage /XYZ @xpos @ypos null]}}
% kein Einstieg: keine Schülerfrage in der Gliederung (Schülerfrage, -Fälle, -Form, -Weg, -Skizze; oder Frage) – Einstieg nur von dort, nie aus der Bank (pruefung.md § 5, 12.10.)
% 2022-bebb-gk-A1.7a-m1 (Original 2022-bebb-gk-A1.7a, Schwere I, rang 58; herausgelöste Fassung: keine)
\begin{pfaufg}{GK~’22\\A1.7a}{1.}{\pfohnetr}
Von den Kunden eines Reisebüros reisen 50 \% gerne in die Region A, 40 \% gerne in die Region B und 15 \% gerne in beide Regionen. Bestimme die Wahrscheinlichkeit, dass ein zufällig ausgewählter A-Kunde auch gerne nach B reist.
\karo{2}
\fusshilfe{\gfuss{1}{0,3}}
\end{pfaufg}
% 2023-bebb-gk-B4.1c-m1 (Original 2023-bebb-gk-B4.1c, Schwere I, rang 59; herausgelöste Fassung: keine)
\begin{pfaufg}{GK~’23\\B4.1c}{2.}{}
Von den Lehrkräften eines Landes arbeiten 30 \% an einem Gymnasium ($G$). 20 \% der Lehrkräfte sind weiblich ($W$) und arbeiten an einem Gymnasium; insgesamt sind 65 \% der Lehrkräfte weiblich. Eine zufällig ausgewählte Lehrkraft ist weiblich. Bestimme die Wahrscheinlichkeit, dass sie an einem Gymnasium arbeitet.
\karo{2}
\fusshilfe{\gfuss{2}{31 \%}}
\end{pfaufg}
% 2022-bebb-gk-B4i-m1 (Original 2022-bebb-gk-B4i, Schwere II, rang 60; herausgelöste Fassung: keine)
\begin{pfaufg}{GK~’22\\B4i}{3.}{}
In einem Paketzentrum haben 12 \% der Pakete das Ziel A und 8 \% das Ziel B; 6 \% aller Pakete sind schwer ($S$). Unter den Paketen mit dem Ziel A ($Z$) sind 10 \% schwer. Untersuche, ob der Anteil der Pakete mit Ziel A unter den schweren Paketen ebenso groß ist wie unter den nicht schweren.
\karo{4}
\fusshilfe{\gfuss{3}{nicht gleich}}
\end{pfaufg}

% ===================== Stufe: bedingte Wahrscheinlichkeit mit Parameter =====================
\gkopfvor{\pfstufekopf{bedingte Wahrscheinlichkeit mit Parameter}{\gpunkte[26 25]{\textcolor{black!80}{$\bullet$}\,\textcolor{black!80}{$\bullet$}\,\textcolor{black!25}{$\circ$}\,\textcolor{black!25}{$\circ$}\,\textcolor{black!25}{$\circ$}}}\special{pdf:dest (stufe-bedingte-wahrscheinlichkeit-mit-parameter) [@thispage /XYZ @xpos @ypos null]}}
% kein Einstieg: keine Schülerfrage in der Gliederung (Schülerfrage, -Fälle, -Form, -Weg, -Skizze; oder Frage) – Einstieg nur von dort, nie aus der Bank (pruefung.md § 5, 12.10.)
% 2025-bebb-gk-B4e-m1 (Original 2025-bebb-gk-B4e, Schwere III, rang 61; herausgelöste Fassung: keine)
\begin{pfaufg}{GK~’25\\B4e}{4.}{}
In einer Stadt haben 31,5 \% der Haushalte mindestens ein Kind. Von diesen Haushalten sind 14,2 \% überbelegt, von den Haushalten ohne Kind sind es 5,5 \%. Zwei Jahre später sind die Anteile ohne Kind unverändert; unter den Haushalten mit Kind ist der Anteil der überbelegten um $x$ gesunken. Die Funktion $f$ mit $f(x) = \tfrac{0{,}315 \cdot (0{,}142 - x)}{0{,}315 \cdot (0{,}142 - x) + 0{,}685 \cdot 0{,}055}$, $0 \le x \le 0{,}142$, ist abgebildet. Bestimme grafisch $x$ mit $f(x) = 0{,}4$ und deute $x$ und $f(x)$ im Sachzusammenhang.
\gzwei{\begin{abiksys}[xmin=0,xmax=0.16,ymin=0,ymax=0.6,xstep=0.02,ystep=0.1,le=50,yle=8]\funktion[0.03]{0.315*(0.142-\x)/(0.315*(0.142-\x)+0.685*0.055)}{G_f}\end{abiksys}}{\karo{4}}
\fusshilfe{\gfuss{4}{0,06 (grafisch, rechnerisch 0,062); 40 \%: Anteil der Haushalte mit Kind unter den überbelegten}}
\end{pfaufg}
% 2026-bb-gk-B4f-m1 (Original 2026-bb-gk-B4f, Schwere III, rang 62; herausgelöste Fassung: keine)
\begin{pfaufg}{GK~’26\\B4f}{5.}{}
In einem Supermarkt sammeln 80 \% der Kundschaft Treuepunkte ($T$). Unter den Sammelnden sind 70 \% weiblich ($w$), unter den übrigen ist der Anteil weiblicher Personen $a$. Ein Jahr später sammeln weiterhin 80 \% Treuepunkte, davon 70 \% weiblich; der Anteil $a$ ist gestiegen. Beurteile die Aussage: Die Wahrscheinlichkeit, dass eine nicht weibliche Person Treuepunkte sammelt, ist größer als vor einem Jahr.
\gzwei{\baumabi{T/{0,8}/{w/{0,7},\overline{w}/},\overline{T}//{w/{a},\overline{w}/}}}{\karo{4}}
\fusshilfe{\gfuss{5}{wahr}}
\end{pfaufg}

% ===================== Stufe: Terme deuten =====================
\gkopfvor{\pfstufekopf{Terme deuten}{\gpunkte[24 22]{\textcolor{black!25}{$\circ$}\,\textcolor{black!25}{$\circ$}\,\textcolor{black!80}{$\bullet$}\,\textcolor{black!25}{$\circ$}\,\textcolor{black!80}{$\bullet$}}}\special{pdf:dest (stufe-terme-deuten) [@thispage /XYZ @xpos @ypos null]}}
% kein Einstieg: keine Schülerfrage in der Gliederung (Schülerfrage, -Fälle, -Form, -Weg, -Skizze; oder Frage) – Einstieg nur von dort, nie aus der Bank (pruefung.md § 5, 12.10.)
% 2022-bebb-gk-B4g-m1 (Original 2022-bebb-gk-B4g, Schwere I, rang 63; herausgelöste Fassung: keine)
\begin{pfaufg}{GK~’22\\B4g}{6.}{}
In einem Paketzentrum haben 12 \% der Pakete das Ziel A und 8 \% das Ziel B; 6 \% aller Pakete sind schwer ($S$). Unter den Paketen mit dem Ziel A ($Z$) sind 10 \% schwer. Gib die Bedeutung der Terme $P_Z(S)$ und $P(S \cap Z)$ im Sachzusammenhang an.
\karo{3}
\fusshilfe{\gfuss{6}{P\_Z(S): Wahrscheinlichkeit, dass ein Paket mit Ziel A schwer ist (10 \%); P(S $\cap$ Z): Wahrscheinlichkeit, dass ein Paket schwer ist und Ziel A hat (0,12 · 0,10 = 1,2 \%)}}
\end{pfaufg}
% 2024-bebb-gk-B4.1a-m1 (Original 2024-bebb-gk-B4.1a, Schwere I, rang 64; herausgelöste Fassung: keine)
\begin{pfaufg}{GK~’24\\B4.1a}{7.}{}
Von der Kundschaft eines Baumarkts sind 70 \% Treuekunden ($T$) und 30 \% Morgenkunden ($M$), die vor 9 Uhr einkaufen. Es gilt $P(\overline{T} \cap M) = 0{,}06$. Deute die Gleichung im Sachzusammenhang.
\karo{2}
\fusshilfe{\gfuss{7}{Die Wahrscheinlichkeit, dass eine zufällig ausgewählte Person Morgenkunde und kein Treuekunde ist, beträgt 6 \%}}
\end{pfaufg}

\pfblattende{10}{26 & 25 & 24 & 23 & 22 & 21 & 20 & 19 & 18 & 17}{B4f & B4e & B4.1a & B4.1c & A1.7a B4g B4i &  & B4.2f &  & B3.2f & B3.1d}%
{\pfausschnitt{Ziehen ohne Zurücklegen \pkt{} Vierfeldertafel}{Bedingte Wahrscheinlichkeit}{Term und Anteil am Baum}}
\end{document}